\scanpage{102}
\mspage{- 10 -}\hfill n\textsuperscript{o} 285

% L1
topologiques $(G_\alpha)$ et la famille d'applications $(j_\alpha\handwritten{\circ}g_\alpha)$, en désignant
% L2
par $g_\alpha$ l'application canonique de $H$ dans $H_\alpha$. Comme $j_\alpha\handwritten{\circ}g_\alpha=f_\alpha\handwritten{\circ}j$,
% L3
le même critère de transitivité prouve que $j$ est un isomorphisme de $H$
% L4
sur le sous-groupe $j(H)$. En outre, $j(H)$ est l'intersection, dans le groupe
% L5
produit $F=\handwritten{\prod}_{\alpha\handwritten{\in}I}G_\alpha$, de $G$ et du sous-groupe $\handwritten{\prod}_{\alpha\handwritten{\in}I}H_\alpha$ d'où la dernière
% L6
assertion du corollaire.

% L7
Lorsque les hypothèses du cor.2 sont satisfaites, on dit que $(H_\alpha)$ est
% L8
un \uline{système projectif de sous-groupes} correspondant au système projectif
% L9--10
de groupes $(G_\alpha)$, et on identifie d'ordinaire $H$ et $j(H)$ au moyen de l'injection $j$.

% L11
Soit $K$ un sous-groupe de $G=\underset{\handwritten{\longleftarrow}}{\lim}\,G_\alpha$, et posons $K_\alpha=f_\alpha(K)$ ; il est
% L12
immédiat que $(K_\alpha)$ est un système projectif de sous-groupes dont la
% L13
\margin{\handwritten{\unsure{2|Z}}}limite projective $H$ s'identifie à un sous-groupe de $G$. Il est clair
% L14
en outre que $H\handwritten{\supset}K$ et on voit aussitôt que $K$ est dense dans $H$
% L15
(cf. prop.23) ; mais en général on aura $H\ne K$ (exerc. voir \handwritten{\S}3,
% L16
cor.1 de la prop.18).

% L17
\uline{Corollaire} 3 - \uline{Les hypothèses et notations étant celles du cor}.1,
% L18--19
\uline{soit} $N_\alpha$ \uline{le noyau de} $g_\alpha$ ; \uline{alors} $(N_\alpha)$ \uline{est un système projectif de sous-groupes distingués}, \uline{dont la limite projective est} \uline{le noyau de} $g=\underset{\handwritten{\longleftarrow}}{\lim}\,g_\alpha$

% L20
La relation $g_\alpha\handwritten{\circ}\varphi_{\alpha\handwritten{\beta}}=\handwritten{\psi}_{\alpha\handwritten{\beta}}\handwritten{\circ}g_{\handwritten{\beta}}$ pour $\alpha\handwritten{\leqslant}\handwritten{\beta}$ montre aussitôt que
% L21
$(N_\alpha)$ est un système projectif. En outre, si $\varphi_\alpha$ (resp. $\handwritten{\psi}_\alpha$) est l'application
% L22
canonique de $A$ dans $A_\alpha$ (resp. de $B$ dans $B_\alpha$), la relation $\handwritten{\psi}_\alpha(g(x))=$
% L23
$=g_\alpha(\varphi_\alpha(x))$ pour $x\handwritten{\in}A$ montre que si $x\handwritten{\in}N=\underset{\handwritten{\longleftarrow}}{\lim}\,N_\alpha$, on a $\handwritten{\psi}_\alpha(g(x))=e$
% L24
pour tout $\alpha\handwritten{\in}I$, donc $g(x)=e$, et réciproquement cette dernière relation
% L25
entraîne $\varphi_\alpha(x)\handwritten{\in}N_\alpha$ pour tout $\alpha\handwritten{\in}I$, donc $x\handwritten{\in}N$. Ce raisonnement montre
% L26
en outre que $N$ est \uline{l'intersection des sous-groupes} $\varphi^{-1}_\alpha(N_\alpha)$ (qui forment
% L27
une base de filtre, comme on le vérifie aussitôt ).
% archive sheet number: 101
